\documentclass[11pt]{article}
\usepackage[T1]{fontenc}
\usepackage{lmodern,microtype,amsmath,amssymb,amsthm,mathtools,mathrsfs}
\usepackage[margin=1in]{geometry}
\usepackage{booktabs,longtable,array,enumitem,tikz-cd}
\usepackage[colorlinks=true,linkcolor=blue,citecolor=blue,urlcolor=blue]{hyperref}
\hypersetup{pdftitle={Local connectedness in recursion categories and essential incompleteness: theorem statements and proof outlines -- version 2},pdfauthor={Florian Lengyel}}
\newtheorem{theorem}{Theorem}[section]
\newtheorem{lemma}[theorem]{Lemma}
\newtheorem{proposition}[theorem]{Proposition}
\newtheorem{corollary}[theorem]{Corollary}
\theoremstyle{definition}\newtheorem{definition}[theorem]{Definition}
\newcommand{\PA}{\mathrm{PA}}
\newcommand{\N}{\mathbb N}
\newcommand{\Set}{\mathbf{Set}}
\newcommand{\Pfn}{\mathbf{Pfn}}
\newcommand{\Tot}{\operatorname{Tot}}
\newcommand{\Split}{\operatorname{Split}}
\newcommand{\Par}{\operatorname{Par}}
\newcommand{\ran}{\operatorname{ran}}
\newcommand{\Th}{\operatorname{Th}}
\newcommand{\Fm}{\mathsf{Fm}}
\newcommand{\Sent}{\operatorname{Sent}}
\newcommand{\Con}{\operatorname{Con}}
\newcommand{\Prov}{\operatorname{Prov}}
\newcommand{\Tr}{\operatorname{Tr}}
\newcommand{\B}{\mathcal B}
\newcommand{\R}{\mathscr R}
\newcommand{\A}{\mathcal A}
\newcommand{\D}{\mathcal D}
\newcommand{\E}{\mathcal E}
\newcommand{\U}{\mathcal U}
\setlength{\emergencystretch}{3em}
\setlist{topsep=4pt,itemsep=3pt}
\title{Local connectedness in recursion categories\\and essential incompleteness\\[1ex]\large Theorem statements and proof outlines}
\author{}
\date{Version 2 --- 1 November 2026}
\begin{document}
\maketitle
\begingroup\small
\tableofcontents
\endgroup
\newpage

\section{Introduction}
Let $F:\R\to\A$ be a range functor from a recursion category.
Theorem~\ref{thm:comparison} constructs a total source morphism whose
interpretation avoids two disjoint target domains containing the images
of the diagonal domains. Its hypothesis is a source lift of the reversed
partial conditional for those target domains. Local connectedness and
weak choice provide such target conditionals
(Proposition~\ref{prop:conditionals}); a Turing morphism in the source
provides the diagonal argument.

Syntax and generated theories are constructed in
$\Tot(\Split\R)$. Pulling the conclusion map of an initial derivation
algebra back to the sentence object gives a map whose range represents
the generated sentence theory. Interpretation transports this range to
the image of the interpreted sentence-conclusion map. A semantic
model gives disjoint theorem and refutation domains, and partial proof
search supplies a conditional lift. For the ambient category $S_{\PA}$,
this proves recursive inseparability of the theorem and refutation sets
of every consistent computably enumerable (c.e.) extension of PA.

The other applications concern computation independent of a binary
representation of Miller's sequence and a recursion category obtained
from cosheaves on a three-point space. The first gives an essentially
incomplete theory of real closed fields under the specified uniform
semidecidability condition. Its inseparable pair acquires a separator
when one representation is fixed. The second gives propositional
theories whose endpoint theories have no complete consistent c.e.
extensions.

Di Paola and Heller proposed categorical formulations of generalized
incompleteness through theorem and refutation sets \cite{DPH}.
Bernardi studied categorical inseparability through precomplete
equivalence relations \cite{Bernardi}. Montagna and Di Paola--Montagna
constructed syntactic recursion categories using provable equality of
partial recursive functions \cite{Montagna,DPM}. The comparison theorem
applies to a source already equipped with a Turing morphism; the
construction of syntax and theories, interpretation, and conditional
lifting provide its logical applications.

\section{Restriction categories, ranges, and recursion}
Composition is written $gf=g\circ f$. A restriction category $\mathcal C$ assigns
to each morphism $f:X\rightharpoonup Y$ a morphism
$\bar f:X\rightharpoonup X$ satisfying
\[
 f\bar f=f,\qquad
 \bar f\bar g=\bar g\bar f,\qquad
 \overline{g\bar f}=\bar g\bar f,\qquad
 \bar h f=f\overline{hf}.
\]
In the middle two equations $g:X\rightharpoonup Z$; in the last,
$h:Y\rightharpoonup Z$. The morphisms $\bar f$ are restriction
idempotents. A morphism is total when $\bar f=1_X$, and the total
morphisms form the category $\Tot(\mathcal C)$
\cite{CockettLack}. A restriction functor preserves the operation
$f\mapsto\bar f$.

A restriction zero $0_{X,Y}$ is absorbing under composition and
satisfies $\overline{0_{X,Y}}=0_{X,X}$; write $0_X=0_{X,X}$.
Restriction idempotents on $X$ are ordered by $a\le b$ when $ba=a$.
Their meet is composition and their greatest element is $1_X$.
Their join, when it exists, is denoted by $a\vee b$. They are
disjoint when $ab=0_X$ and complementary when also $a\vee b=1_X$.
A disjoint pair $a_0,a_1$ is inseparable if no complementary pair
$c,c'$ satisfies $a_0\le c$ and $a_1\le c'$.
For parallel partial morphisms, the restriction order is
$f\le g$ when $f=g\bar f$; joins of morphisms refer to this order.
Parallel morphisms $f,g$ are compatible when $f\bar g=g\bar f$.

A restriction idempotent $r$ on $Y$ receives
$f:X\rightharpoonup Y$ when $rf=f$. The range $\ran(f)$ is the
least such restriction idempotent. Ranges satisfy the agreement law
\[
 gf=hf\quad\Longrightarrow\quad
 g\ran(f)=h\ran(f)
 \qquad(g,h:Y\rightharpoonup Z).
\]
In $\Pfn$, the category of sets and partial functions, restriction
and range are the partial identities on the domain and image.

A binary restriction product $X\times Y$ has total projections
$\pi_1,\pi_2$ and, for $f:Z\rightharpoonup X$ and
$g:Z\rightharpoonup Y$, a unique pairing with
\[
 \pi_1\langle f,g\rangle=f\bar g,\qquad
 \pi_2\langle f,g\rangle=g\bar f,\qquad
 \overline{\langle f,g\rangle}=\bar f\bar g.
\]
Its total diagonal is $\Delta_X=\langle1_X,1_X\rangle$.
A binary restriction coproduct is a categorical coproduct with total
injections $i_0:X\to X+Y$ and $i_1:Y\to X+Y$. Write $[f,g]$ for
copairing and $\nabla_X=[1_X,1_X]:X+X\to X$ for the codiagonal.
Ranges and restriction coproducts give binary joins of restriction
idempotents by $a\vee b=\ran[a,b]$.

For $f:X\rightharpoonup Y$ and a restriction idempotent $e$ on $Y$,
$f^*e=\overline{ef}$. A range functor
$F:\mathcal C\to\mathcal C'$ between categories with ranges,
restriction zeros, and binary restriction coproducts preserves these
structures and restrictions. It consequently
preserves joins because $a\vee b=\ran[a,b]$. Write
\[
 \chi_{X,Y}:FX+FY\xrightarrow{\cong}F(X+Y)
\]
for its coproduct comparison. The definition does not include preservation
of restriction products.

A restriction idempotent $e$ on $X$ splits when $e=mr$ for a total
$m:E\to X$ and a partial $r:X\rightharpoonup E$ with $rm=1_E$.
A restriction category is split when every restriction idempotent
splits. The objects of $\Split\mathcal C$ are pairs $(X,e)$ with
$e$ a restriction idempotent. A morphism
$f:(X,e)\rightharpoonup(Y,d)$ is a morphism of $\mathcal C$
satisfying $f=dfe$. The identity of $(X,e)$ is $e$, and $f$ is
total precisely when $\bar f=e$. Ranges extend to this category:
$df=f$ implies $\ran(f)\le d$, so the same representative is a
range on $(Y,d)$. The inclusion $X\mapsto(X,1_X)$ is fully
faithful. In particular, a morphism $U\rightharpoonup U+U$ between
original objects constructed in the splitting belongs to the source.

A restriction idempotent $a$ is connected if every disjoint
decomposition $a=b\vee c$ has $b=0$ or $c=0$. A category is
locally connected when its connected restriction idempotents are
jointly epimorphic: $fa=ga$ for every connected restriction
idempotent $a$ on the common source implies $f=g$.
A partial monomorphism $p:X\rightharpoonup Y$ satisfies
$pu=pv\Rightarrow\bar p u=\bar p v$ for all
$u,v:W\rightharpoonup X$. A section of $p$ is a partial morphism
$s:Y\rightharpoonup X$ with $ps=\bar s$ and $psp=p$.
Weak choice means that every partial monomorphism has a section.

A restriction category with binary restriction products and zeros
is prodominical when $f\times0=0$ for all morphisms of the required
types, and is isotypical when its objects are pairwise isomorphic.
A Turing morphism $\tau:U\times U\rightharpoonup U$ has the
following property: for every object $X$ and every
$f:X\times U\rightharpoonup U$, there is a total parameter
$r:X\to U$ with $f=\tau(r\times1_U)$. A recursion category is
a prodominical isotypical category with binary restriction products
and a Turing morphism. These structures express partial computation
in the restriction-category framework
\cite{CockettLack,CockettHofstra}.

\section{Diagonal arguments and lifted partial conditionals}

\begin{theorem}[Comparison theorem]\label{thm:comparison}
Let $\R$ be a recursion category with ranges and binary restriction
coproducts, and let $\A$ be a restriction category with ranges,
restriction zeros, and binary restriction coproducts. Let
$F:\R\to\A$ be a range functor. Choose a Turing morphism
$\tau:U\times U\rightharpoonup U$ and an isomorphism
$b:U\to U+U$. Put
\[
 h=b\tau\Delta_U,\qquad
 a_0=h^*(1_U+0_U),\qquad a_1=h^*(0_U+1_U).
\]
Let $P,N$ be disjoint restriction idempotents on $FU$ such that
$F(a_0)\le P$ and $F(a_1)\le N$. Suppose a partial morphism
$s:FU\rightharpoonup FU+FU$ satisfies
\[
 s^*(1_{FU}+0_{FU})=N,\qquad
 s^*(0_{FU}+1_{FU})=P,
\]
and suppose there is $\ell:U\rightharpoonup U+U$ with
$F\ell=\chi_{U,U}s$.
Then there is a total $g:U\to U$ such that
\[
 (Fg)^*P=(Fg)^*N=0_{FU}.
\]
If $F(1_U)\ne0$ and inverse image along $Fg$ preserves the join
$P\vee N$, then $P\vee N\ne1_{FU}$.
\end{theorem}
\begin{proof}
Put $u=b^{-1}\ell:U\rightharpoonup U$. Apply the Turing property
to $u\pi_2:U\times U\rightharpoonup U$, obtaining a total
$g:U\to U$ with $\tau(g\times1_U)=u\pi_2$.
Precompose with $\langle1_U,g\rangle:U\to U\times U$ to obtain
$\tau\Delta_Ug=ug$, hence $hg=\ell g$.

Write $E=Fg$ and $H=\chi_{U,U}^{-1}Fh$. Then $HE=sE$.
Preservation of restrictions and the coproduct comparison give
\[
 E^*F(a_0)=E^*N,\qquad E^*F(a_1)=E^*P.
\]
The two inequalities in the statement give $E^*N\le E^*P$
and $E^*P\le E^*N$. Their meet is
$E^*(PN)=0$, so both are zero. Totality of $E$ and preservation
of the indicated join would give
$1_{FU}=E^*1_{FU}=E^*(P\vee N)=0$ if $P\vee N=1_{FU}$.
\end{proof}

\begin{corollary}[Separators whose conditionals lift]\label{cor:separator}
Under the categorical assumptions of Theorem~\ref{thm:comparison},
assume $F(1_U)\ne0$ and that total inverse images in $\A$ preserve
binary joins. A complementary pair $P,N$ with
$F(a_0)\le P$ and $F(a_1)\le N$ cannot have a reversed conditional
admitting a source lift as in that theorem.
\end{corollary}
\begin{proof}
Apply the theorem and use $P\vee N=1_{FU}$.
\end{proof}
Every disjoint pair in $\Pfn$ has a set-theoretic separator.
Corollary~\ref{cor:separator} constrains source lifts of the
associated reversed conditional.

\begin{lemma}[Connected domains and coproducts]\label{lem:connected}
Suppose $\A$ has the restriction structure above and inverse images
preserve binary joins. If $a$ is a connected restriction idempotent
on $X$ and $f:X\to Y+Z$ is total, then
$(1_Y+0_Z)fa=fa$ or $(0_Y+1_Z)fa=fa$.
\end{lemma}
\begin{proof}
Pull the disjoint summand decomposition back along $f$ and meet
with $a$. Connectedness makes one piece zero; the other is $a$.
The restriction identity gives the asserted equation.
\end{proof}

\begin{proposition}[Conditionals in a locally connected category]\label{prop:conditionals}
Let $\A$ be split and locally connected, with ranges, restriction
zeros, and binary restriction coproducts. Assume weak choice and
preservation of binary joins by inverse images. For disjoint
restriction idempotents $P,N$ on $X$, the partial morphism
\[
 p=\nabla_X(N+P):X+X\rightharpoonup X
\]
has a unique section $s:X\rightharpoonup X+X$. It satisfies
\[
 \bar s=P\vee N,\qquad
 s^*(1_X+0_X)=N,\qquad s^*(0_X+1_X)=P.
\]
\end{proposition}
\begin{proof}
For total $u,v:Y\to X+X$, Lemma~\ref{lem:connected} places each
of their restrictions to a connected domain in one summand.
The equation $pu=pv$ then gives $\bar p u=\bar p v$ on every
connected domain, and hence everywhere by local connectedness.
For arbitrary partial $u,v$, splitting the common domain of $pu,pv$
reduces to this total case. Thus $p$ is a partial monomorphism.
Weak choice supplies a section. The section equations
give $\bar s=\ran(p)=P\vee N$; partial monicity gives uniqueness.
The identities $sP=i_1P$ and $sN=i_0N$, together with this domain
equation, give the two summand inverse images.
\end{proof}

\begin{corollary}[Inseparability in the source]\label{cor:source}
Let $F:\R\to\A$, $\tau$, $b$, and $a_0,a_1$ be as in
Theorem~\ref{thm:comparison}, with $\A$ as in
Proposition~\ref{prop:conditionals} and $F(1_U)\ne0$.
If the reversed conditional associated with $Fc,Fc'$ lifts for every
complementary source pair $c,c'$ on $U$, then $a_0,a_1$ are inseparable
in $\R$.
\end{corollary}
\begin{proof}
A source separator is sent to a complementary target pair, since
a range functor preserves joins, meets, and zero. The assumed lift
contradicts Corollary~\ref{cor:separator}.
\end{proof}
Appendix C gives a faithful range functor with a separable diagonal
pair for which the required conditional has no source lift.

\begin{proposition}[Disjoint joins give source lifts]\label{prop:joins}
If parallel partial morphisms with disjoint domains have joins in
the source, then for complementary $c,c'$ on $X$ the morphism
\[
 \ell=(i_1c)\vee(i_0c'):X\to X+X
\]
is total and satisfies $\ell c=i_1c$ and $\ell c'=i_0c'$.
If $F:\R\to\A$ is a range functor and $\A$ satisfies
Proposition~\ref{prop:conditionals}, then
$\chi_{X,X}^{-1}F\ell$ is the target section of that proposition.
\end{proposition}
\begin{proof}
The join equations give the restrictions on $c,c'$; their union
is the full domain. After the coproduct identification, both maps agree on
$Fc,Fc'$. Their copairing has full range, so range agreement gives
equality.
\end{proof}

\section{Internal theories and semantic consistency}
Let $\D=\Tot(\Split\R)$ satisfy the hypotheses used in Appendix A.
Assume the semantic target $\A$ is split, and choose an extension
$\widehat F:\Split\R\to\A$ preserving ranges and the restriction
structure. This extension is the composite of the canonical functor
to $\Split\A$ and the functor given by chosen splittings in $\A$.
The identity splitting on original objects gives
$\widehat F|_{\R}=F$.
Let $\Phi\in\D$ be the formula object and $j:S\hookrightarrow\Phi$
its sentence subobject. For an axiom subobject $A\hookrightarrow S$,
let $c_A:P_A\to\Phi$ be the initial derivation algebra's conclusion.
Form the pullback in $\D$:
\[
\begin{tikzcd}
 P_{A,S}\arrow[r]\arrow[d,"c_{A,S}"']&P_A\arrow[d,"c_A"]\\
 S\arrow[r,hook,"j"']&\Phi.
\end{tikzcd}
\]
Put $t_A=\ran(c_{A,S})$ in $\Split\R$ and let
$\nu:S\to S$ be total syntactic negation. Appendix A constructs
these objects and morphisms from the syntax and inference-rule data.

\begin{proposition}[Interpretation of the theorem object]\label{prop:image}
Let $F:\R\to\A$ and $\widehat F$ be as specified above. Then
\[
 \widehat F(t_A)=\ran(\widehat F(c_{A,S})).
\]
If the target is presented as $\Par(\E,\operatorname{Mono})$ for
a regular category $\E$, a factorization in $\E$ of the total
conclusion map through a validity monomorphism $V\hookrightarrow
\widehat F S$ implies that the image of the internal theorem object
lies in $V$.
\end{proposition}
\begin{proof}
The first equation is preservation of ranges. Apply the range/image
lemma of Appendix A to the conclusion map in the total category;
factorization through a monomorphism bounds its regular image.
\end{proof}
The image of $\widehat F(c_{A,S})$ is the interpreted sentence-theorem
object. It includes the conclusions of every code in the interpreted
proof object, including nonstandard proof codes in arithmetic models.

\begin{theorem}[Incompleteness from semantic consistency]\label{thm:logical}
Let the source data be as in Theorem~\ref{thm:comparison}. Suppose the
target has total inverse images preserving binary joins, with
$F(1_U)\ne0$. Let $\theta:U\to S$ be total in $\Split\R$.
Put $t=\widehat F(t_A)$ and define restriction idempotents on $FU$ by
\[
 P=(\widehat F\theta)^*t,\qquad
 N=(\widehat F\theta)^*(\widehat F\nu)^*t.
\]
Assume $t(\widehat F\nu)^*t=0$, $F(a_0)\le P$, and
$F(a_1)\le N$. If the reversed partial conditional for $P,N$
has a source lift, there is a total $g:U\to U$ such that
\[
 (\widehat F(\theta g))^*t=0,\qquad
 (\widehat F(\theta g))^*(\widehat F\nu)^*t=0.
\]
Consequently $t\vee(\widehat F\nu)^*t\ne1_{\widehat F S}$.
\end{theorem}
\begin{proof}
Target consistency makes $P,N$ disjoint. Apply
Theorem~\ref{thm:comparison}. The two displayed equations are
its conclusion after composing inverse images. If the two theorem
domains covered $\widehat F S$, their pullbacks along the total
$\widehat F(\theta g)$ would cover $FU$, contradicting
$F(1_U)\ne0$.
\end{proof}

\begin{corollary}[The source-consistency criterion]\label{cor:source-logical}
Let a restriction category have zeros and binary joins preserved by
total inverse images. Let $a,b$ be inseparable on $U$, let
$\theta:U\to S$ and $\nu:S\to S$ be total, and let $t_0$ be a
restriction idempotent on $S$ satisfying
$a\le\theta^*t_0$ and $b\le\theta^*\nu^*t_0$.
Every restriction idempotent $t$ on $S$ with $t\ge t_0$
and $t\nu^*t=0$ has inseparable theorem and refutation domains
$t,\nu^*t$; in particular, $t\vee\nu^*t\ne1_S$.
\end{corollary}
\begin{proof}
A complementary separator of $t,\nu^*t$ pulls back to one of
$a,b$. The same argument applies when $t,\nu^*t$ themselves are
assumed complementary.
\end{proof}

\section{Syntactic recursion categories and their internal syntax}
Throughout this section, $B\supseteq\PA$ denotes the ambient
arithmetic theory and $A$ denotes a represented axiom presentation.
Fix the program enumeration with its finite computation traces.
Write $G(e,x,y)$ for its $\Sigma_1$ computation-graph predicate,
defined by the existence of a finite computation trace, and put
$G_e(x,y)=G(e,x,y)$. Choose primitive-recursive numerical codings
of products and tagged coproducts with PA-provable decoding identities.

\begin{lemma}[Arithmetic computation laws]\label{lem:computation}
For the fixed program enumeration and numerical codings, the graph
predicates are $\Sigma_1$ and
PA-provably functional. PA proves the uniform graph identities for
composition, identities, projections, pairing, coproduct injections,
copairing, domains, and the chosen universal evaluator. In particular,
there are a universal index $u$ and primitive-recursive functions
$\operatorname{s},\operatorname{comp}:\N\times\N\to\N$ for which
PA proves the universal closures of
\[
\begin{aligned}
G_u(\langle e,x\rangle,y)&\leftrightarrow G(e,x,y),\\
G(\operatorname{s}(e,a),x,y)&\leftrightarrow G(e,\langle a,x\rangle,y),\\
G(\operatorname{comp}(f,e),x,z)&\leftrightarrow
 \exists y\,[G(e,x,y)\land G(f,y,z)].
\end{aligned}
\]
For each fixed $e$, the function $a\mapsto\operatorname{s}(e,a)$
has a program whose graph and totality PA proves.
\end{lemma}
\begin{proof}
The computation laws are proved by the finite-trace transformations
used in Montagna's formalized computation argument \cite{Montagna}.
Composition concatenates two traces and records the intermediate
value. The structural programs implement their numerical codings.
The evaluator simulates the indexed program, and the parameter
program gives its uniform factorization. All equalities are universal
graph equalities, not separate proofs for individual numerals.
\end{proof}
The objects of $S_B$ are generated from the numerical object
$\omega$ by the chosen products and coproducts. A morphism is a
program index with the indicated source and target, modulo
$B$-provable uniform graph equality. Composition is represented
by the function $\operatorname{comp}$.

\begin{lemma}[Functional search]\label{lem:search}
For a $\Sigma_1$ relation $\Gamma(x,y)$ there is a partial recursive
search program with domain $\exists y\,\Gamma(x,y)$; if $B$ proves
functionality of $\Gamma$, it proves uniform equality of that
program's graph with $\Gamma$.
\end{lemma}
\begin{proof}
Search increasing finite bounds for a witness to the primitive
recursive matrix. The trace verification proves existence exactly
on the existential domain. Under functionality, any discovered
witness has the unique permitted output.
\end{proof}

\begin{lemma}[The total splitting]\label{lem:total}
For consistent $B\supseteq\PA$, the category
$\B_B=\Tot(\Split S_B)$ is equivalent to the category of
$\Sigma_1$ carrier predicates and $B$-provably total functional
$\Sigma_1$ relations, modulo $B$-provable uniform graph equality.
It is regular with disjoint pullback-stable finite coproducts.
Every monomorphism is an admissible domain monomorphism, and
\[
 \Par(\B_B,\operatorname{Mono})\simeq\Split S_B.
\]
\end{lemma}
\begin{proof}
Use Lemma~\ref{lem:search} to identify functional relations with
program morphisms. Products and equalizers use conjunction and
numerical pairing; tagged unions give extensive coproducts. Images
are existential graph predicates. A graph whose surjectivity is proved
in $B$ has a witness-search section, so it is the coequalizer of its
kernel pair and remains so
under pullback. For a monomorphism, its kernel-pair projections are
equal; this proves injectivity uniformly in $B$. Searching its
transposed graph gives its partial inverse. These constructions give
the asserted partial-map equivalence.
\end{proof}

\begin{lemma}[Finite membership conditions]\label{lem:finite}
PA proves that a coded finite family of $\Sigma_1$ membership
conditions holds exactly when a coded finite family of witnesses
exists. Consequently the finite validity tests used below are
$\Sigma_1$ predicates with the required PA-provable equivalences.
\end{lemma}
\begin{proof}
Induction on the coded list length appends a witness at each step.
Conversely each component of a witness list verifies its membership
condition. This is induction in PA, so it covers internally finite
lists in every model of $B$.
\end{proof}

\begin{proposition}[Syntax and lists with parameters]\label{prop:syntax}
Fix bounded finite arity and variable and signature objects of
$\B_B$. The numerical term and formula tree objects realize the
appendix's initial syntax algebra, stable under every parameter
object. The numerical list object realizes the corresponding
parameterized initial list algebra.
\end{proposition}
\begin{proof}
Use canonical trees and lists as carriers. Numerical equality makes
every represented variable carrier decidable in $\B_B$.
For an arbitrary algebra
over a parameter object, an evaluation trace records the intermediate
values, the algebra-operation witnesses, and the common parameter.
Lemma~\ref{lem:finite} makes validity $\Sigma_1$. Induction in $B$
proves existence and uniqueness of evaluation, and a second induction
proves uniform graph equality with every algebra morphism. This gives
initiality in each slice.
\end{proof}

\begin{proposition}[Substitution realization]\label{prop:substitution}
The numerical free-variable, free-for-substitution, term-substitution,
and formula-substitution operations realize the categorical folds and
equations specified in Appendix A in $\B_B$.
\end{proposition}
\begin{proof}
The numerical constructor equations, including the quantifier case
and the free-for condition, are provable in PA. Parameterized initiality identifies
these operations with the appendix's folds; uniqueness gives the
identities as $B$-provable graph equalities.
\end{proof}

\begin{lemma}[Initial derivations]\label{lem:proof-initial}
Every bounded-arity rule presentation by objects and total morphisms
of $\B_B$ has an initial derivation algebra in the slice over its
formula object.
\end{lemma}
\begin{proof}
Use trees labeled by rule instances with ordered children.
Finite membership conditions and the matching of premises to child
conclusions are $\Sigma_1$. Evaluation into any rule algebra is
$\Sigma_1$; $B$-induction proves its existence, functionality,
algebra equations, and uniqueness. The constructors are labeled by
rule instances; membership witnesses occur in the $\Sigma_1$ validity test.
\end{proof}

\begin{theorem}[Internally generated theories]\label{thm:internal}
For consistent $B\supseteq\PA$, the bounded-arity first-order
structure of Appendix A is realized in $\B_B$. Every axiom
subobject $A\hookrightarrow S$ has an initial derivation algebra.
Its conclusion image is the least subobject closed under the
specified rules, and its sentence pullback is the generated sentence
theory. Let $e_A$ be the program whose domain is exactly the
valid-derivation predicate $P_A$ and whose output is the conclusion.
The partial identity on the numerical codes of the sentence theory,
extended by zero outside the sentence carrier, is
\[
 t_A^{\omega}=[\Sent(z)\land\exists p\,G_{e_A}(p,z)]_B:
 \omega\rightharpoonup\omega.
\]
Its restriction to the sentence splitting is
$t_A:S\rightharpoonup S$.
\end{theorem}
\begin{proof}
Apply Lemmas~\ref{lem:total} and \ref{lem:proof-initial},
Propositions~\ref{prop:syntax} and \ref{prop:substitution}, and the
least-theory theorem of Appendix A. Functional search realizes the
sentence-conclusion pullback. The range/image lemma identifies its
program range with the regular image in the total category.
\end{proof}
For $B=\PA$, a c.e. axiom presentation $A$ determines the axiom
subobject of $S$ with carrier predicate
$\Sent(z)\land\operatorname{Ax}_A(z)$. Its inclusion is represented
by the identity on sentence codes.

\begin{proposition}[Numerical realization]\label{prop:rec-realization}
The category $\Tot(\Split\mathrm{Rec})$ is equivalent to the category
of c.e. numerical carriers and partial recursive functions total on
their source carriers. It realizes the same bounded-arity first-order
syntax, substitution, lists, and generated-theory construction.
\end{proposition}
\begin{proof}
Finite limits use pairing and equality, coproducts use tags, and
images are c.e. by dovetailing. Evaluation of ordinary finite trees
gives initiality and uniqueness by induction. Finite nonlogical
axiom membership tests are semidecidable. The range/image lemma gives
the theorem partial identity.
\end{proof}

\paragraph{Computational inseparability and completeness.}
For $B=\Th(\N)$, interpreting each program in $\N$ gives an
isomorphism $S_{\Th(\N)}\cong\mathrm{Rec}$ with the chosen common
object presentation. Two indices are equal in the source exactly
when their uniform graph-equivalence sentence belongs to
$\Th(\N)$, which is equivalent to equality of their standard
partial functions. Thus $S_{\Th(\N)}$ has inseparable computational
pairs although $\Th(\N)$ is complete. This example distinguishes
the existence of an inseparable pair from its representation by
the theorem and refutation domains of a specified internal theory.

\section{Model interpretation, soundness, and ultrapowers}

\begin{proposition}[Interpretation in one model]\label{prop:model}
Let $B\supseteq\PA$ be consistent and let $M\models B$.
Interpreting the program graph formulas defines a range functor
$F_M:S_B\to\Pfn$ preserving restriction products. Its extension
to splittings, with the subset splittings in $\Pfn$, restricts to
\[
 G_M:\B_B\longrightarrow\Set
\]
preserving finite limits, finite coproducts, and regular images.
\end{proposition}
\begin{proof}
Interpret the uniform computation identities of
Lemma~\ref{lem:computation} in $M$. $B$-provable graph equality
makes the assignment independent of indices. In the predicate
presentation, a carrier $C$ becomes
$C(M)=\{x\in M:M\models C(x)\}$. The finite-limit formulas,
tagged coproducts, and existential image formulas become the
corresponding set constructions. Internal traces and witnesses range
over all of $M$, including its nonstandard elements.
\end{proof}

\begin{proposition}[Source complementarity and inseparability]\label{prop:st-pair}
For complementary restriction idempotents $c,c'$ on $X$ in $S_B$,
there is a total $s:X\to X+X$ with $sc=i_1c$ and $sc'=i_0c'$.
For any $M\models B$, its image is the reversed conditional for
$F_Mc,F_Mc'$. Consequently the diagonal pair is inseparable in
$S_B$ for every consistent $B\supseteq\PA$.
\end{proposition}
\begin{proof}
Search for a witness to either $\Sigma_1$ domain predicate.
$B$-proved disjointness and coverage give the uniform graph identity
and totality in the source. Interpretation gives the required
conditional. The target $\Pfn$ satisfies
Proposition~\ref{prop:conditionals}, and $F_M(1_U)\ne0$.
Apply Corollary~\ref{cor:source}. A model exists by completeness.
\end{proof}
Compatibility with extension of the ambient theory is described
in Appendix D.

\begin{proposition}[Soundness from one semantic model]\label{prop:soundness}
Fix the ambient category $\B_{\PA}$, a c.e. first-order arithmetic
axiom presentation $A$, its sentence-conclusion map
$c_{A,S}:P_{A,S}\to S$, and a model $K\models A$.
Let $\nu:S\to S$ be syntactic negation. For a standard sentence
code $s$, let $\varphi_s$ denote the sentence it codes.
Define a subobject in $\Set$ by
\[
 V_K=\{s\in G_{\N}S:K\models\varphi_s\},\qquad
 v_K:V_K\hookrightarrow G_{\N}S.
\]
There is a function $e_K:G_{\N}P_{A,S}\to V_K$ with
$v_Ke_K=G_{\N}c_{A,S}$. Moreover
$V_K\cap(G_{\N}\nu)^{-1}V_K=\varnothing$, and $A$ is consistent.
\end{proposition}
\begin{proof}
In the standard numerical interpretation, every proof code represents
a standard finite derivation in the fixed calculus. Its nonlogical
axioms hold in $K$; the logical axiom schemes are valid; modus ponens
and generalization preserve validity. Induct on the entire derivation,
taking validity of an open formula to mean truth under every
assignment in $K$. At a sentence conclusion this is
$K\models\varphi_s$, which gives the factorization.
A sentence and its negation cannot both
hold in $K$. Since every standard finite derivation occurs in this carrier, a
pair of contradictory proofs would contradict that disjointness.
\end{proof}
The functor $G_{\N}$ interprets the numerical code objects.
The structure $K$ determines the validity subobject $V_K$ of
the sentence-code carrier.

The ultrapower construction \cite[Theorem 3.1 and Proposition 4.2]{Keisler}
gives the following functor on sets.
\begin{lemma}[The ultrapower functor]\label{lem:ultrafunctor}
Work in a set-theoretic metatheory with choice. Let $\U$ be a proper
ultrafilter on a nonempty set $I$. For a set $X$ and two families
$(x_i),(y_i)\in X^I$, define
\[
 (x_i)\sim_{\U}(y_i)
 \quad\Longleftrightarrow\quad
 \{i\in I:x_i=y_i\}\in\U.
\]
The assignments
\[
 \mathscr U(X)=X^I/{\sim_{\U}},\qquad
 \mathscr U(f)([x_i])=[f(x_i)]
\]
for sets $X$ and functions $f:X\to Y$ define a faithful functor
$\mathscr U:\Set\to\Set$ preserving
finite limits, finite coproducts, surjections, and regular images.
Here $X^I$ is the set of functions $I\to X$ in $\Set$.
\end{lemma}
\begin{proof}
Composition is coordinatewise; constant families detect inequality
of functions. The empty set and a singleton are fixed by $\mathscr U$.
Pairing proves finite-product preservation. An
equalizer condition holds on a $\U$-large set; replace the other
coordinates by a fixed element of that equalizer to obtain a
representative in its ultrapower. If the equalizer is empty, both
sets are empty. For a binary coproduct, the ultrafilter selects one
of the two tag sets, and the same replacement gives a representative
entirely in that summand. Together with the empty-coproduct case,
this proves finite-coproduct preservation.
Surjections have sections by choice, and their
section equations are preserved. Surjection--injection factorizations
therefore remain image factorizations.
\end{proof}

\begin{theorem}[Ultrapower comparison]\label{thm:ultra}
Let $\U$ be a proper ultrafilter on a nonempty set $I$, and let
$\mathscr U$ be the functor of Lemma~\ref{lem:ultrafunctor}.
Let $B\supseteq\PA$ satisfy $\N\models B$, and put
$M=\N^I/\U$. There is a natural isomorphism
\[
 \eta:\mathscr U G_{\N}\xrightarrow{\cong}G_M
 \quad\text{on }\B_B.
\]
It identifies the ultrapowers of the standard sentence-conclusion
maps and their images with their interpretations in $M$.
\end{theorem}
\begin{proof}
For a carrier predicate $C$, \L{}o\'s's theorem identifies
$C(M)$ with classes having $C(\N)$-valued coordinates on a large
set. Replace other coordinates by a fixed member of $C(\N)$;
the empty case gives empty sets on both sides. For a functional
graph $\Gamma$, \L{}o\'s identifies its interpreted graph with the
ultrapower graph. This proves naturality. Finite-limit and image
preservation give the conclusion about sentence pullbacks and
theorem images \cite{Keisler}.
\end{proof}

\begin{proposition}[Failure of initial-algebra preservation]\label{prop:ultra-initiality}
Let $I=\N$ and let $\U$ be a nonprincipal ultrafilter on $\N$.
The $1+X$-algebra obtained by applying $\mathscr U$ to
$(\N,0,\mathrm{succ})$ is not initial in $\Set$.
\end{proposition}
\begin{proof}
The element $[i]$ of $\N^{\N}/\U$ differs from every numeral,
since a singleton does not belong to $\U$. The unique algebra
morphism from $(\N,0,\mathrm{succ})$ maps each natural number
to its numeral class. This morphism is not surjective. If its target
were also initial, it would be an isomorphism.
\end{proof}

\begin{theorem}[Soundness for the full interpreted proof object]\label{thm:ultrasound}
Under Proposition~\ref{prop:soundness}, let $I,\U,\mathscr U$
be as in Lemma~\ref{lem:ultrafunctor}, put $M=\N^I/\U$,
and use Theorem~\ref{thm:ultra} with ambient $B=\PA$.
There is a factorization
\[
 G_MP_{A,S}\xrightarrow{e_{K,\U}}\mathscr U(V_K)
 \xrightarrow{v_{K,\U}}G_MS,
 \qquad v_{K,\U}e_{K,\U}=G_Mc_{A,S}.
\]
Its validity subobject and its inverse image under $G_M\nu$ are
disjoint. Hence the interpreted internal theorem and refutation
domains are disjoint, including their nonstandard codes.
\end{theorem}
\begin{proof}
Apply $\mathscr U$ to $v_Ke_K=G_{\N}c_{A,S}$ and use the
components of $\eta$. Explicitly,
\[
 e_{K,\U}=\mathscr U(e_K)\eta_{P_{A,S}}^{-1},\qquad
 v_{K,\U}=\eta_S\mathscr U(v_K).
\]
Naturality gives the equation, and preservation of pullbacks and
the empty set gives disjointness. Proposition~\ref{prop:image}
identifies the interpreted internal theorem object with this
conclusion image.
\end{proof}

\paragraph{The assignment sort.}
Consider the many-sorted structure consisting of the numerical
code sort $\N$, the semantic domain $K$, the assignment sort
$K^{\N}$, and the evaluation, assignment-modification, and satisfaction
relations. Its ultrapower has code sort $\N^I/\U$, semantic sort
$K^I/\U$, and assignment sort $\mathscr U(K^{\N})$. The transferred
evaluation map has type
\[
 \mathscr U(K^{\N})\times\mathscr U(\N)\longrightarrow\mathscr U(K).
\]
The compositional clauses for connectives, quantification, and
substitution transfer by \L{}o\'s's theorem. Satisfaction is a named
relation of this semantic expansion. For $K=\N\models A$, this
is the ultrapower of the standard arithmetic structure expanded
by its satisfaction relation.

\begin{corollary}[Consistency under ultrapowers]\label{cor:consistency}
Let $I,\U,\mathscr U$ be as in Lemma~\ref{lem:ultrafunctor}.
Let $A$ be any consistent c.e. arithmetic theory represented in
$\B_{\PA}$, and put $M=\N^I/\U$. Its interpreted theorem and
refutation domains are disjoint. No assumption $\N\models A$
is required.
\end{corollary}
\begin{proof}
The standard conclusion image is the theorem set of $A$,
so its negation intersection is empty by consistency. Apply the
finite-limit and image preservation of $\mathscr U$ and the
comparison isomorphism. Equivalently, transfer the true arithmetic
sentence asserting that no sentence has both an $A$-proof and an
$A$-refutation.
\end{proof}

\begin{proposition}[Limits of model interpretation]\label{prop:limits}
The section-lifting hypotheses do not imply soundness for every model
interpretation. Moreover, taking an ultrapower of a fixed model does
not improve faithfulness of its graph interpretation.
\end{proposition}
\begin{proof}
Assuming consistency of PA, choose
$L\models\PA+\neg\Con(\PA)$, with the proof predicate for the
chosen calculus. The interpretation $G_L$ contains a proof code whose
conclusion is $0=1$. Its conclusion map cannot factor through any
validity subobject excluding that sentence, although the source
complementary-conditionals lift as in Proposition~\ref{prop:st-pair}.
Every ultrapower of $L$ retains this sentence. The example is also
used in the partial-combinatory-algebra literature \cite[\S7]{Terwijn}.

For the faithfulness assertion, equality of two interpreted program
graphs is expressed by their uniform graph-equality sentence. A
model and its ultrapower satisfy the same such sentences. Thus
their functors identify exactly the same pairs of source morphisms.
\end{proof}

\section{Arithmetic: theorem sets, refutations, and essential incompleteness}
Fix the arithmetic signature, the internal syntax of
Theorem~\ref{thm:internal}, and its sentence object $S\in\B_{\PA}$.
For every c.e. axiom presentation $A$, the standard interpretation
of its generated-theory subobject is the set of consequences of $A$
with standard finite derivations: its carrier predicate recognizes exactly the
finite derivation trees of the chosen calculus.
Write $\Prov_A(z)$ for the $\Sigma_1$ sentence-conclusion image
predicate supplied by Theorem~\ref{thm:internal}, including the
sentence predicate and the chosen axiom presentation.

\begin{lemma}[Representation of the diagonal domains]\label{lem:reduction}
Choose the numerical diagonal domains $D_0,D_1$ obtained by
interpreting $a_0,a_1$ in $\N$. There are $\Sigma_1$ formulas
$\alpha(x),\beta(x)$ and a total morphism in $\Split S_{\PA}$
\[
 \theta:\omega\to S,\qquad
 (\widehat F_{\N}\theta)(n)=\ulcorner\alpha(\bar n)\urcorner,
\]
such that $\alpha,\beta$ define $D_0,D_1$ in $\N$,
$\PA\vdash\forall x\,\neg(\alpha(x)\land\beta(x))$, and
\[
 n\in D_0\Rightarrow\PA\vdash\alpha(\bar n),\qquad
 n\in D_1\Rightarrow\PA\vdash\neg\alpha(\bar n).
\]
Consequently, for every c.e. $A\supseteq\PA$, the target reductions
\[
 F_{\N}(a_0)\le(\widehat F_{\N}\theta)^*\widehat F_{\N}(t_A),\qquad
 F_{\N}(a_1)\le(\widehat F_{\N}\theta)^*
                  (\widehat F_{\N}\nu)^*\widehat F_{\N}(t_A)
\]
hold, where $\widehat F_{\N}$ is the chosen extension to the
restriction splitting. They also hold for $F_M,\widehat F_M$ in
every $M\equiv\N$.
\end{lemma}
\begin{proof}
Use finite computation traces ending with the two output tags.
Determinism proves disjointness in PA. A true trace can be verified
by a finite PA proof; for the second tag, combine that proof with
the disjointness theorem. Numeral construction and substitution are
the morphisms realized in \S5. Extension includes these PA theorems
among the $A$-theorems. The resulting uniform arithmetic implications
are true in $\N$, so they hold in $M\equiv\N$.
\end{proof}

Let $\theta_{\mathrm{num}},\nu_{\mathrm{num}}:\N\to\N$ be the
chosen primitive recursive code functions realizing $\theta$ and
$\nu$ on their specified carriers. Thus
$\theta_{\mathrm{num}}(n)=\ulcorner\alpha(\bar n)\urcorner$,
and $\nu_{\mathrm{num}}$ forms the negation code. Inside arithmetic
formulas, these function symbols abbreviate their PA-provably
functional and total graph definitions.

\begin{lemma}[A partial proof-search lift]\label{lem:proof-search}
Let $A$ be a c.e. axiom presentation. There is a partial morphism
$\ell_A:\omega\rightharpoonup\omega+\omega$ in $S_{\PA}$
searching for a proof of the sentence coded by
$\theta_{\mathrm{num}}(x)$ or of its negation and returning
the reversed tag. In any model $M\models\PA$ in which these two
proof predicates are disjoint, its interpretation is the partial
conditional for their two domains.
\end{lemma}
\begin{proof}
Let $C(x)=\Prov_A(\theta_{\mathrm{num}}(x))$ and
$C'(x)=\Prov_A(\nu_{\mathrm{num}}(\theta_{\mathrm{num}}(x)))$.
Enumerate witnesses to the two
$\Sigma_1$ predicates in increasing finite bounds, with a fixed
priority if both are first found at the same stage. PA proves that
the search halts exactly on $C\lor C'$ and that every returned
value satisfies
\[
 (C(x)\land z=\iota_1(x))\lor
 (C'(x)\land z=\iota_0(x)).
\]
Here $\iota_j(x)=2x+j$ realizes the injection $i_j$.
In a model where $C,C'$ are disjoint, this relation is functional,
so the interpreted search has exactly the displayed graph.
\end{proof}

\begin{theorem}[Incompleteness from the internal theory object]\label{thm:arithmetic-incomplete}
For every consistent c.e. extension $A\supseteq\PA$, the theory
object constructed in $\B_{\PA}$ has an incomplete standard
interpretation. There is a total source morphism
$g:\omega\to\omega$ whose standard interpretation
$\gamma=F_{\N}(g):\N\to\N$ satisfies: no number $\gamma(n)$ belongs to either
of the inverse images of its theorem and refutation sets under
the numerical substitution function $\theta_{\mathrm{num}}$. In particular,
$\alpha(\overline{\gamma(0)})$ is undecidable in $A$.
\end{theorem}
\begin{proof}
Consistency gives disjoint target domains in the standard interpretation.
For an ultrapower interpretation, use
Corollary~\ref{cor:consistency}.
Lemma~\ref{lem:reduction} supplies the target inequalities and
Lemma~\ref{lem:proof-search} supplies the partial lift.
Theorem~\ref{thm:logical} gives the total $g$ and the two zero
inverse images. For the individual ordinary sentence, use the
standard interpretation and evaluate at $0$.
\end{proof}

\begin{theorem}[Recursive inseparability and essential incompleteness]\label{thm:arithmetic-insep}
For every consistent c.e. extension $A\supseteq\PA$, its theorem
and refutation sets of sentence codes are recursively inseparable.
Every consistent c.e. extension of $A$ is therefore incomplete.
\end{theorem}
\begin{proof}
Suppose a recursive set separates the theorem and refutation sets.
Pull it back along the recursive numeral-substitution function
$\theta_{\mathrm{num}}$.
Its two complementary domains contain $D_0,D_1$ by
Lemma~\ref{lem:reduction}. Its reversed conditional is a total
recursive function. Choose a program index for it; that index is
a partial morphism of $S_{\PA}$, regardless of whether PA proves
its totality. Its standard interpretation is exactly the target
conditional. Corollary~\ref{cor:separator} gives a contradiction.
The same proof applies to every consistent c.e. extension. A
complete consistent c.e. theory would have complementary c.e.
theorem and refutation sets, whose simultaneous enumeration decides
them, contradicting recursive inseparability.
\end{proof}
In $\Pfn$, disjoint subsets have set-theoretic separators.
The proof applies to recursive separators because their conditional
functions are represented by program indices in $S_{\PA}$.

\paragraph{Extension maps.}
For c.e. axiom presentations $A\subseteq A'$ within the fixed
ambient $\B_{\PA}$, the generated-theory construction gives
$\Th_A\le\Th_{A'}$ whenever the chosen axiom inclusion is a
morphism of the source. A presentation of $A'$ containing the chosen
presentation of $A$ as an explicit tagged component gives a
PA-verifiable inclusion. A common PA component verifies only the
inclusions from PA into the two presentations. For the
external essential-incompleteness statement, ordinary inclusion of
theories and target inclusion suffice. The identity on the standard sentence
carrier gives a covariant map between the corresponding disjoint
theorem/refutation pairs. These maps vary the axiom presentation
within a fixed ambient category. The quotient functor
$S_B\to S_{B'}$ instead changes the ambient theory.

\paragraph{Comparison with the numerical proof.}
The standard interpretation of the internal objects in $\B_{\PA}$
gives the c.e. proof and theorem carriers of
Proposition~\ref{prop:rec-realization}. The arithmetic argument uses the
source Turing morphism and the interpreted source proof-search
morphism.

\section{Miller's sequence and a theory of real closed fields}
Computation uniform over all binary representations of a real
sequence determines a recursion category. The comparison theorem
produces an inseparable pair in this category, and Miller's index
property reduces that pair to axioms of a theory of real closed fields.

Fix a recursive pairing $\langle-,-\rangle:\N\times\N\to\N$
and an acceptable oracle-program enumeration $(\Phi_e^p)_e$.
Put $Q=[0,1]^{\N}$ and define the binary representation map
\[
 \rho:\{0,1\}^{\N}\to Q,\qquad
 \rho(p)(n)=\sum_{j\ge0}p(\langle n,j\rangle)2^{-j-1}.
\]
Write $P_z=\rho^{-1}(\{z\})$. Each $P_z$ is nonempty and compact.
Fix a recursive rational enumeration $(q_j)_j$. A program $d$
computes a real $r$ from $p$ when $\Phi_d^p(k)$ is defined for
every $k$ and $|q_{\Phi_d^p(k)}-r|<2^{-k}$. Choose a Miller
sequence $a\in Q$ with the index property: if $d$ computes the
same real $r\in[0,1]$ from every $p\in P_a$, then $a(d)=r$.
The index property is the strengthened conclusion following
Miller's Theorem~6.3; the binary-representation formulation is
given by Bauer--Hanson
\cite[Theorem~6.3 and following paragraph]{Miller}
\cite[\S2.1 and Theorem~2.2]{BauerHanson}.

\begin{lemma}[Unanimous computation]\label{lem:miller-unanimity}
There is one oracle program which, from $p\in P_z$ and numerical
inputs $e,n$, computes the partial function
$\tau_z:\N\times\N\rightharpoonup\N$ defined by
\[
 \tau_z(e,n)=m\quad\Longleftrightarrow\quad
 (\forall q\in P_z)\;\Phi_e^q(n)\mathord\downarrow=m.
\]
Its value and domain depend only on $z$, not on the selected
representation $p$.
\end{lemma}
\begin{proof}
Finite halting computations enumerate the open set
$H_{e,n,m}=\{q:\Phi_e^q(n)\mathord\downarrow=m\}$ by cylinders.
From $p\in P_z$, enumerate cylinders $[\sigma]$ for which
$z\notin\rho([\sigma])$: the image is a finite union of closed
dyadic boxes, and strict inequalities are semidecidable from $p$.
These cylinders enumerate the complement of $P_z$. For each $m$,
search for a finite cylinder cover of Cantor space using this
complement and $H_{e,n,m}$. The finite-cover test is decidable;
compactness makes success equivalent to $P_z\subseteq H_{e,n,m}$.
Dovetail these searches over $m$. Nonemptiness of $P_z$ ensures
uniqueness of the returned value, and the success condition depends
only on $z$.
\end{proof}

Let $\E_a$ have the objects generated from $\N$ by binary
products and tagged coproducts, with fixed recursive codings.
A morphism is a numerical partial function computed by one oracle
program with the same domain and values for every $p\in P_a$;
morphism equality is equality of partial functions.

\begin{proposition}[The uniform recursion category]\label{prop:miller-category}
The category $\E_a$ is a locally connected recursion category with
ranges, binary restriction coproducts, finite compatible joins,
and weak choice. The morphism
$\tau_a:\N\times\N\rightharpoonup\N$ is a Turing morphism.
\end{proposition}
\begin{proof}
The structural programs preserve uniform domains and values.
Preimage search computes ranges and, for injective morphisms,
partial inverse sections; the returned value is independent of
the representation. Compatible computations can be dovetailed
to obtain their join. Lemma~\ref{lem:miller-unanimity} supplies
the evaluator, and the oracle parameter theorem supplies its total
parameter maps. Singleton restrictions detect equality of partial
functions and prove local connectedness. The numerical codings
give isotypicality, and products with zero are zero.
\end{proof}

Choose the recursive coproduct isomorphism that tags a natural
number by its parity. The resulting diagonal domains are
\[
 A_i=\{e:\tau_a(e,e)\mathord\downarrow=m
                 \text{ for some }m\equiv i\pmod 2\},
 \qquad i=0,1.
\]
Let
$L_a=\{n:a(n)<1/2\}$ and $R_a=\{n:a(n)>1/2\}$.

\begin{lemma}[Reduction to the cut-index sets]\label{lem:miller-cuts}
There is an ordinary total recursive function $h:\N\to\N$ such
that $h^{-1}(L_a)=A_0$ and $h^{-1}(R_a)=A_1$.
\end{lemma}
\begin{proof}
For each $e$, the sets
$U_e^i=\{z:\tau_z(e,e)\downarrow=m\text{ for some }m\equiv i\pmod2\}$
are disjoint effectively open subsets of $Q$. Finite halting-cylinder
covers express each as a union of complements of computably
specified finite unions of closed dyadic boxes. Weighted sums of
distances to these finite unions give uniformly computable continuous functions
$g_{e,i}:Q\to[0,1]$ that are positive exactly on $U_e^i$.
The function
$r_e(z)=\tfrac12+(g_{e,1}(z)-g_{e,0}(z))/4$
is uniformly computable from any binary representation of $z$.
The parameter theorem gives an ordinary total recursive index
function $h$ for these real-valued computations. Miller's index
property gives $a(h(e))=r_e(a)$, proving both inverse-image
equalities \cite{Miller}.
\end{proof}

\begin{corollary}[Inseparability and fixed-representation separation]\label{cor:miller-cuts}
The partial identities on $L_a,R_a$ are inseparable in $\E_a$.
For each fixed binary representation $p$ of $a$, a $p$-decidable
set separates the two underlying numerical sets.
\end{corollary}
\begin{proof}
Use source conditionals and Corollary~\ref{cor:source} for the
diagonal pair, then pull a hypothetical cut separator back along
$h$. The first binary digit of the fixed representation gives the
$p$-decidable separator. Thus the faithful range functor
$J_p:\E_a\to\mathrm{Rec}^{p}$ sends this pair to a separable
pair. If it preserved the chosen Turing morphism, the target
diagonal theorem would contradict that separator. Hence it does not.
\end{proof}

\begin{proposition}[First-order realization in the total splitting]\label{prop:miller-realization}
The category $\Tot(\Split\E_a)$ realizes the bounded-arity
first-order syntax, substitution, lists, and generated-theory
construction of Appendix A.
\end{proposition}
\begin{proof}
Carriers are uniformly semidecidable
from all representations. Pairing, tags, equality, and uniform
search give the total categorical structure. Evaluate finite syntax
and derivation trees with the common operations; their domains and
values remain independent of the representation. Induction proves
initiality, and ranges give the generated theories.
\end{proof}

\begin{theorem}[Essential incompleteness in the uniform computability class]\label{thm:miller-incomplete}
Let $T_a$ consist of a recursive axiomatization of real closed
ordered fields in the language with constants $(c_n)_{n\in\N}$,
together with the axioms
\[
 q<c_n\quad(q<a(n)),\qquad c_n<q\quad(a(n)<q),
 \qquad q\in\mathbb Q, n\in\N.
\]
Every consistent extension of $T_a$ whose axiom set has a partial
identity belonging to $\E_a$ has inseparable theorem and refutation
domains in $\Split\E_a$, and is incomplete.
\end{theorem}
\begin{proof}
The real field with $c_n=a(n)$ gives consistency of $T_a$.
The realized proof object gives the theorem domains for each
extension. The map $n\mapsto\ulcorner c_n<1/2\urcorner$ reduces
$L_a,R_a$ to the theorem and refutation domains. Consistency
is source disjointness in this extensional category.
Corollary~\ref{cor:source-logical} gives the conclusion.
\end{proof}

\section{Three-point cosheaves and propositional theories}
The three-point cosheaf construction gives a recursion category
not restriction equivalent to $\mathrm{Rec}$. Its two endpoint
interpretations yield recursively inseparable theorem and refutation
sets for internally generated propositional theories.
Let $\mathcal C_{\mathrm{rec}}$ be the category of c.e. numerical
carriers and partial recursive functions total on their source
carriers, as in Proposition~\ref{prop:rec-realization}.
Use the fork diagram category $\E=\mathcal C_{\mathrm{rec}}^{J}$,
where $J$ has objects $0,1,2$ and nonidentity arrows $0\to1,0\to2$,
its partial-map category
$\A_X=\Par(\E,\operatorname{Mono})$, and the recursion subcategory
$\R_X$ constructed in Appendix B. Its splitting is equivalent
to $\A_X$. Identify $\E$ with $\Tot(\A_X)$ through the canonical
equivalence. The chosen object
$V\in\E$ has components and structure maps
\[
 \begin{gathered}
 V_0=\N\times(\N\times(\N\times\N)),\qquad
 V_1=V_2=\N\times\N,\\
 r_1^V(k,z,u,v)=(k,u),\qquad
 r_2^V(k,z,u,v)=(k,v).
 \end{gathered}
\]
Appendix B constructs a Turing morphism
$\tau:V\times V\rightharpoonup V$ and an isomorphism
$\zeta:V\to V+V$ whose component tags are determined by the
parity of the first coordinate. Put
\[
 a=(\zeta\tau\Delta_V)^*(1_V+0_V),\qquad
 b=(\zeta\tau\Delta_V)^*(0_V+1_V).
\]
Write $D_a,D_b\hookrightarrow V$ for their c.e. domain subfunctors
and put $A_j=(D_a)_j$, $B_j=(D_b)_j$, for $j=0,1,2$.
The full inclusion $\R_X\hookrightarrow\A_X$ lifts the
conditionals supplied by Proposition~\ref{prop:conditionals}.
Corollary~\ref{cor:source} makes $a,b$ inseparable; fullness retains
this conclusion in $\A_X$.
The object $V$ does not have decidable equality, so this application
uses the propositional construction of Appendix A.

Put $\mathcal K=\Par(\mathcal C_{\mathrm{rec}},\operatorname{Mono})$.
Component evaluation gives functors
$E_j:\A_X\to\mathcal K$ for $j=0,1,2$; on a span it takes
the span of its $j$th components.

\begin{proposition}[Endpoint interpretation]\label{prop:fork-endpoints}
The functors $E_j$ preserve restrictions, zeros, ranges, binary
restriction products and coproducts, and finite joins of restriction
idempotents. For $i=1,2$, the morphism
\[
 E_i(\tau):(\N\times\N)\times(\N\times\N)
                  \rightharpoonup\N\times\N
\]
is a Turing morphism in $\mathcal K$, and $A_i,B_i$ are recursively
inseparable. The morphism $E_0(\tau)$ is not Turing in $\mathcal K$.
\end{proposition}
\begin{proof}
Component evaluation preserves the indicated operations because
they are computed componentwise. With a recursive bijection
$\operatorname{encode}_2:\N\times\N\to\N$ and its inverse
$\operatorname{decode}_2$, the endpoint evaluator has the formula
\[
 (E_i\tau)((k,p),(\ell,x))
 =\operatorname{decode}_2\varphi_p(\operatorname{encode}_2(\ell,x)).
\]
The parameter theorem gives the required total parameter maps.
Let $\mathcal K_2$ be the full subcategory of
$\mathcal K$ generated from $\N\times\N$ by the chosen binary
products and coproducts. Its objects are recursively isomorphic to
$\N\times\N$, and the displayed evaluator and its parameter maps
make it a recursion category. The underlying-partial-function
functor $\mathcal K_2\to\Pfn$ is a range functor.
Complementary c.e. domains have total recursive conditionals, so
Corollary~\ref{cor:source} gives the inseparability of the endpoint
diagonal pair. A recursive separator would be a complementary
restriction pair in $\mathcal K_2$.
For the zeroth component, consider
$f:V_0\to V_0$, $f(k,z,u,v)=(0,0,z,0)$: naturality would make
$r_1^Vf(k,z,u,v)$
independent of $z$ if $f$ were a specialization of $E_0(\tau)$.
\end{proof}

Let $S\in\E$ have finite propositional formula trees as its
components, with atom labels in $V_j$ and structure maps induced
by relabeling along $r_i^V$. Its constructors are morphisms
\[
 \theta:V\to S,\qquad \nu:S\to S,
 \qquad {\Rightarrow}:S\times S\to S.
\]
The logical axioms are the instances of the schemes
\[
 \begin{gathered}
 \varphi\Rightarrow(\psi\Rightarrow\varphi),\\
 (\varphi\Rightarrow(\psi\Rightarrow\chi))
 \Rightarrow((\varphi\Rightarrow\psi)\Rightarrow(\varphi\Rightarrow\chi)),\\
 (\nu\psi\Rightarrow\nu\varphi)\Rightarrow(\varphi\Rightarrow\psi).
 \end{gathered}
\]
Together with the positive and negative atomic axiom morphism
$[\theta|_{D_a},\nu\theta|_{D_b}]:D_a+D_b\to S$ and modus ponens,
these schemes specify the derivation rules. The finite derivation
trees give an object $\mathrm{Prf}\in\E$
and a conclusion morphism $c:\mathrm{Prf}\to S$.

\begin{proposition}[Propositional syntax and derivations]\label{prop:fork-initiality}
The constructor algebra of $S$ is initial for the endofunctor
$Z\mapsto V+Z+(Z\times Z)$ on $\E$. The derivation algebra
$(\mathrm{Prf},c)$ is initial for the rule endofunctor on $\E/S$
determined by these axioms and modus ponens.
\end{proposition}
\begin{proof}
Evaluate finite trees at each component. Induction proves that
evaluation commutes with the two fork structure maps, hence is a
morphism of $\E$; componentwise uniqueness proves initiality.
The same argument applies to derivations with their conclusion
labels.
\end{proof}

View $c$ as a total morphism in $\A_X$ and put $t_0=\ran(c)$,
a restriction idempotent on $S$. Let $D_{t_0}\hookrightarrow S$
be its domain subfunctor.

\begin{theorem}[Generated component theories]\label{thm:fork-generated}
The subfunctor $D_{t_0}$ is the least c.e. subfunctor of $S$
containing the logical and specified positive and negative atomic
axioms and closed under modus ponens. Its component theories are
consistent, and
\[
 a\le\theta^*t_0,\qquad b\le\theta^*\nu^*t_0,
 \qquad t_0\nu^*t_0=0_S.
\]
\end{theorem}
\begin{proof}
Apply Appendix A's least-theory and range/image results to the
initial derivation algebra. One-node proofs give the two
inequalities. For any alleged pair of contradictory finite proofs,
make precisely their positive atomic labels true. Disjointness of
the positive and negative atomic axioms makes every leaf valid;
induction contradicts opposite conclusions. This gives componentwise
consistency and hence source disjointness in $\A_X$.
\end{proof}

\begin{theorem}[Compatible extensions and endpoint incompleteness]\label{thm:fork-extensions}
Every compatible family of consistent c.e. deductively closed
extensions of the component theories has inseparable theorem and
refutation subobjects in $\A_X$. At each endpoint, their inverse
images under the atom map are disjoint recursively inseparable c.e.
sets. Extension inclusions give identity-on-carrier morphisms
between these pairs, preserving both distinguished subsets, and
these assignments form a covariant functor. Each endpoint theory
is incomplete.
\end{theorem}
\begin{proof}
Use Corollary~\ref{cor:source-logical} in $\A_X$ for the first
assertion. At endpoint $i$, the two c.e. inverse images contain
$A_i,B_i$. A recursive separator would separate that endpoint
diagonal pair. Completeness would make the inverse images
complementary c.e. sets, hence recursive. Monotonicity of both
sets under extension gives the asserted identity-on-carrier
morphisms of pairs.
\end{proof}

\begin{corollary}[Essential incompleteness of the endpoint theories]\label{cor:fork-essential}
Every consistent c.e. extension of either endpoint base theory is
incomplete.
\end{corollary}
\begin{proof}
For either endpoint $i$, the theorem and refutation inverse images
of any consistent c.e. extension are disjoint c.e. sets containing
$A_i,B_i$. A complete extension would make these inverse images
complementary and hence recursive, contradicting the recursive
inseparability of $A_i,B_i$.
\end{proof}

\appendix
\section{Syntax and deductively closed theories}
The categorical construction has six parts.
\begin{enumerate}
\item \emph{Categorical hypotheses and notation.} Regularity,
  disjoint stable finite coproducts where needed, explicitly
  parenthesized finite powers, $\mathbf2=1+1$, Boolean morphisms
  defined by copairings, and conditional decompositions.
\item \emph{First-order syntax.} The ambient category $\D$, the
  variable object $V\in\D$ with decidable equality, bounded-arity
  signature objects, the paired syntax endofunctor
  $K(X,Y)=(K_1(X,Y),K_2(X,Y))$, parameterized initiality, and lists.
\item \emph{Substitution.} Free-variable and free-for constructions,
  term and formula substitution, quantifier cases, and their
  universal-property proofs.
\item \emph{Inference rules and closed subobjects.} Rule objects,
  premise and conclusion maps, the slice endofunctor, its initial
  derivation algebra, and the image argument.
\item \emph{First-order theories.} The specified logical axiom
  schemes, generalization, modus ponens, sentence pullback, and
  generated-theory closure operator.
\item \emph{Ranges and regular images; propositional instance.}
  The precise comparison with partial-map ranges and the
  propositional specialization needed by \S9.
\end{enumerate}

Let $\D$ be a regular category with disjoint pullback-stable finite
coproducts. Write $\mathbf2=1+1$ with injections
$\mathbf t,\mathbf f:1\to\mathbf2$.
Fix $V\in\D$ with decidable equality: there is a morphism
$\epsilon:V\times V\to\mathbf2$ whose pullback along $\mathbf t$
is the diagonal $\Delta_V:V\hookrightarrow V\times V$.
For a fixed integer $d\ge0$, let $\Sigma_k,\Pi_k\in\D$ be
the objects of function and relation symbols of arity $k$, for
$0\le k\le d$. Finite powers mean the chosen products
$X^0=1$ and $X^{n+1}=X^n\times X$.
Define $K:\D^2\to\D^2$ by
$K(X,Y)=(K_1(X,Y),K_2(X,Y))$, where
\[
 \begin{aligned}
 K_1(X,Y)&=V+\coprod_{k=0}^d\Sigma_k\times X^k,\\
 K_2(X,Y)&=(X\times X)+\coprod_{k=0}^d\Pi_k\times X^k
                    +Y+(Y\times Y)+(V\times Y).
 \end{aligned}
\]
The action on morphisms uses products and copairings on the
displayed summands. The formula constructors are equality,
relations, negation, implication, and universal quantification.

\begin{definition}[Categorical syntax]
A syntax object is an initial algebra
$\iota:K(\mathsf{Tm},\Fm)\to(\mathsf{Tm},\Fm)$ stable under
parameters. There is also an initial list algebra
$[\mathsf{nil},\mathsf{cons}]:1+V\times W\to W$, stable under
parameters. Here stability means that, for every $P\in\D$,
pullback along $P\to1$ gives initial algebras for the corresponding
syntax endofunctor on $(\D/P)^2$ and list endofunctor on $\D/P$,
with all constant objects pulled back along the same morphism.
\end{definition}
Initial algebras specify abstract syntax by its constructors and
recursion property \cite{GTWW}. The parameterized list recursion
used here is the list-object structure of \cite{CockettLists}.

The free-variable fold $\mathsf{FV}:\Fm\to W$ concatenates the
free-variable lists of arguments and deletes the bound variable at
a quantifier. Deletion is defined by list recursion using $\epsilon$.

\begin{definition}[Sentence object]
The sentence object is the pullback
\[
\begin{tikzcd}
 S\arrow[r,hook,"j"]\arrow[d] & \Fm\arrow[d,"\mathsf{FV}"]\\
 1\arrow[r,"\mathsf{nil}"'] & W.
\end{tikzcd}
\]
\end{definition}

\begin{proposition}[Categorical substitution]
Under the appendix's parameterized syntax and list initiality and
decidable-variable hypotheses, the specified free-variable,
free-for, and substitution operations exist as morphisms of $\D$
and satisfy their constructor equations, including quantifiers.
\end{proposition}
\begin{proof}
An enlarged algebra keeps the original syntax component together
with the recursively computed output.
Parameterized initiality gives the fold; its projection to syntax
is the identity by uniqueness. Conditional copairings use the
declared decidable predicates. In the quantifier case, deletion
and the free-for predicate determine whether substitution is
admissible; the same fold yields its constructor equations.
\end{proof}

For the rule construction, regularity and finite coproducts suffice.
Fix $\Phi\in\D$ and, for $0\le n\le d$, morphisms
$p_n:R_n\to\Phi^n$ and $q_n:R_n\to\Phi$ giving the premises
and conclusion of a rule instance. Nullary rules are axioms.
For $q:Q\to\Phi$, put
\[
 B_n(Q,q)=R_n\times_{\Phi^n}Q^n,\qquad
 H(Q,q)=\left(\coprod_{n=0}^d B_n(Q,q),
                   [q_n\operatorname{pr}_1]_{n=0}^d\right).
\]
Pullback uniqueness defines the action of the endofunctor
$H:\D/\Phi\to\D/\Phi$ on morphisms.

\begin{definition}[Deductively closed subobject]
A subobject $m:T\hookrightarrow\Phi$ is closed under the rules
if $(T,m)$ admits an $H$-algebra structure in $\D/\Phi$.
Equivalently, there is a morphism
$\alpha:\coprod_n B_n(T,m)\to T$ satisfying
$m\alpha=[q_n\operatorname{pr}_1]_n$.
For a logical calculus, these are its deductively closed formula
subobjects.
\end{definition}
The algebra structure is unique because $m$ is monic. This is
the categorical form of finitary closure by conclusions of finite
proofs \cite{Aczel}.

\begin{theorem}[Least rule-closed subobject]
Let $\D$ be regular with finite coproducts, let $\Phi\in\D$,
and fix rule morphisms $p_n:R_n\to\Phi^n$, $q_n:R_n\to\Phi$
for $0\le n\le d$.
If the induced endofunctor $H:\D/\Phi\to\D/\Phi$ has an
initial algebra with conclusion $c:P\to\Phi$, then the regular
image of $c$ is the least subobject closed under these rules.
\end{theorem}
\begin{proof}
The rule endofunctor preserves regular epimorphisms: its summands
are pullbacks of finite powers, followed by a finite coproduct.
Factor $c=me$. The regular epimorphism $H(e)$ and the monomorphism
$m$ give a unique diagonal defining an algebra structure on the
image. Initiality maps $P$ into every other closed subobject;
the image factorization gives the required inclusion. The image
algebra construction is an instance of lifting the regular
epimorphism--monomorphism factorization to algebras for an
endofunctor preserving regular epimorphisms \cite{Wissmann}.
\end{proof}

\begin{theorem}[Generated first-order theory]
With the specified categorical first-order syntax and an axiom
subobject of the sentence object, an initial derivation algebra
gives the least formula subobject containing all axioms and closed
under the rules. Its sentence pullback is the least generated
sentence theory. When these derivation algebras exist for every
axiom subobject, the construction is an extensive, monotone,
idempotent closure operator on sentence subobjects.
\end{theorem}
\begin{proof}
Apply the preceding theorem to the specified rule objects.
Regular images commute with the sentence pullback. Leastness gives
extensivity and monotonicity; the generated theory is already
closed, which gives idempotence.
\end{proof}

\begin{lemma}[Regular images and ranges]
Let $\mathcal C$ be a split restriction category with ranges,
and suppose $\Tot(\mathcal C)$ is regular. Let a total morphism
$f:X\to Y$ have regular-image factorization $f=me$ in
$\Tot(\mathcal C)$. If $m:I\to Y$ has a partial inverse
$r:Y\rightharpoonup I$ with $rm=1_I$ and $mr=\bar r$, then
$\ran(f)=mr$. The corresponding assertion for a partial map
applies to its total part on its split domain.
\end{lemma}
\begin{proof}
The image idempotent receives the map. Every receiving restriction
idempotent supplies a monomorphism through which the total part
factors, so regular-image leastness makes the image idempotent
smaller. This is the range universal property; the categorical
relation between ranges and regular images is developed in
\cite{RangeII}.
\end{proof}

\section{The three-point cosheaf recursion category}
Let $X=\{a,b,c\}$ have open sets
$\varnothing,\{c\},\{a,c\},\{b,c\},X$. The cosheaf construction
on this space has the following ten steps. The categories
$\mathcal C_{\mathrm{rec}}$, $\E=\mathcal C_{\mathrm{rec}}^J$,
and $\A_X=\Par(\E,\operatorname{Mono})$, and the object $V$,
are as in \S9. The nonempty-word object $V^\#$ is formed
componentwise, with structure maps acting on each letter.
\begin{longtable}{@{}p{.13\textwidth}p{.30\textwidth}p{.50\textwidth}@{}}
\toprule
Result & Statement & Proof\\\midrule\endhead
B.1 & Set-valued cosheaves on $X$ are equivalent to fork diagrams. & Restriction to the three proper nonempty opens gives the fork. Its pushout gives the cosheaf value on $X$, providing a quasi-inverse.\\
B.2 & The category $\E$ is regular and extensive, with componentwise monomorphisms and images. & Finite limits, tagged coproducts, and c.e. images are computed componentwise; the induced morphisms satisfy the fork equations.\\
B.3 & The category $\A_X$ has ranges, weak choice, and the required restriction products and coproducts. & Image factorizations from B.2 give ranges. Partial inverse injections give sections, and the span construction supplies the restriction structure.\\
B.4 & The category $\A_X$ is locally connected. & A singleton at an endpoint, or a singleton at component $0$ together with its two images, gives a connected restriction. These restrictions detect partial morphisms.\\
B.5 & Every object admits a restriction embedding into $V$. & Recursive component codes give a natural injection whose inverse on its c.e. image is a partial morphism.\\
B.6 & There are isomorphisms $V\cong V\times V$, $V\cong V+V$, and $V\cong V^\#$. & Pairing, tagged parity, and finite-word coding commute with both fork maps.\\
B.7 & The evaluator $\tau:V\times V\rightharpoonup V$ is a Turing morphism. & Ordinary numerical parameterization gives the endpoint maps; the middle component checks the two endpoint agreements. The resulting total parameter map satisfies all three component equations.\\
B.8 & The full subcategory $\R_X$ generated from $V$ is a recursion category and $\Split\R_X\simeq\A_X$. & B.3 and B.6--B.7 give the recursion structure. The restriction embeddings in B.5 give essential surjectivity after splitting.\\
B.9 & The Karoubi envelope of $\R_X$ is equivalent to $\A_X$. & Every idempotent of $\A_X$ splits through its image; B.8 then gives the equivalence.\\
B.10 & The category $\R_X$ is not restriction equivalent to $\mathrm{Rec}$, and $V$ has no decidable equality. & The terminal object has five subobjects in $\E$ and two in $\mathcal C_{\mathrm{rec}}$. Distinct middle-component elements with the same endpoint images prevent a complemented diagonal on $V$.\\
\bottomrule
\end{longtable}

\section{Why lifting cannot be omitted}
\begin{proposition}[A range comparison without the required lifts]
There is a recursion-category example with a faithful range functor
preserving restriction products and reflecting zero maps for which
the source-inseparability conclusion fails when the section-lifting
hypothesis is removed.
\end{proposition}
\begin{proof}
Let $K=\{e:\varphi_e(e)\downarrow\}$ be the ordinary halting set,
and let $(W_a^K)_a$ enumerate the $K$-c.e. subsets of $\N$.
The category $\mathcal C_K$ has the numerical objects generated
from $\N$ by products and coproducts. Its morphisms have $K$-c.e.
domains and ordinary partial recursive extensions. Composition,
restrictions, products, coproducts, and ranges preserve this class:
their domains and images are $K$-c.e., and their value maps have
the indicated recursive extensions. Its total morphisms are exactly
the ordinary total recursive functions. The underlying-function
functor $F:\mathcal C_K\to\Pfn$ is faithful and preserves the
specified structures and reflects zero.

Using a recursive pairing with projections $\pi_0,\pi_1$, define
\[
 \tau_K(\langle e,a\rangle,x)\simeq
 \begin{cases}
 \varphi_e(x),&x\in W_a^K,\\
 \text{undefined},&x\notin W_a^K.
 \end{cases}
\]
Specializing the extension program and the oracle-domain program
gives an ordinary total recursive parameter map. Thus $\tau_K$
is Turing, and the numerical codings make $\mathcal C_K$ isotypical.

For the parity coproduct isomorphism, the two diagonal domains
are separated by $A$ and $A^c$, where
\[
 A=\{q:\varphi_{\pi_0(q)}(q)\downarrow
                          \text{ with an even value}\}.
\]
Both sets are $K$-c.e., so they define complementary source
restrictions. The set $A$ is not recursive: choose a total recursive
$h$ such that $\varphi_{h(n)}$ ignores its input, simulates
$\varphi_n(n)$, and returns $0$ on halting. For a fixed $a_*$,
the map $n\mapsto\langle h(n),a_*\rangle$ reduces $K$ to $A$.
The reversed conditional
for this pair is total and decides $A$ by its output tag. It is
therefore not recursive and cannot be the image of a source morphism.
The target $\Pfn$ has all the other hypotheses of
Corollary~\ref{cor:source}, so this is the omitted lift.
The two disjoint tagged restrictions on $A$ and $A^c$ have no
common upper bound either: such an upper bound would be this same
nonrecursive total conditional.
\end{proof}

\section{Faithful families of models and ambient theory extension}
The family of all countable model interpretations detects every
failure of provable uniform graph equality.

\begin{theorem}[Faithful family of graph interpretations]
For consistent $B\supseteq\PA$, let $\mathcal M_B$ be the set
of arithmetic models of $B$ with underlying set $\N$, and put
\[
 \D_B=\prod_{M\in\mathcal M_B}\Pfn,
 \qquad F_B=(F_M)_M:S_B\to\D_B.
\]
Then $F_B$ is faithful and preserves restrictions, zeros, ranges,
binary restriction products and coproducts, and finite joins of
restriction idempotents. The target is split and locally connected, satisfies
weak choice and inverse-image preservation of joins, and
$F_B(1_\omega)\ne0$.
\end{theorem}
\begin{proof}
Preservation is coordinatewise. A graph equality not provable in
$B$ fails in a model, and completeness with downward
L\"owenheim--Skolem supplies a countable such model, transported
to the carrier $\N$. Singletons at individual coordinates detect
equality and give local connectedness. Partial inverse injections
give weak choice. Consistency makes the model family nonempty.
\end{proof}

\begin{proposition}[Ambient theory extension]
For consistent $B\subseteq B'$, the quotient functor
$Q_{B,B'}:S_B\to S_{B'}$ and restriction of model families
$\pi_{B,B'}:\D_B\to\D_{B'}$ give the commuting square
\[
\begin{tikzcd}
 S_B\arrow[r,"Q_{B,B'}"]\arrow[d,"F_B"'] & S_{B'}\arrow[d,"F_{B'}"]\\
 \D_B\arrow[r,"\pi_{B,B'}"'] & \D_{B'}.
\end{tikzcd}
\]
The quotient functor preserves the defining operations of the
chosen diagonal pair and sends it to the inseparable pair over
$B'$. These assignments compose for successive theory extensions.
\end{proposition}
\begin{proof}
Provability is monotone, so the quotient on indices is well defined.
Each coordinate of either composite interprets the same graph.
The quotient is surjective on hom-sets; a preimage of any target
morphism supplies its Turing factorization. Apply
Proposition~\ref{prop:st-pair} over $B'$.
\end{proof}
The image $F_B(\tau_B)$ is not a Turing morphism in $\D_B$.
At any fixed coordinate $M$, its specializations are indexed by
the countable set $|M|$, whereas there are uncountably many partial
functions $|M|\rightharpoonup|M|$. A partial function outside this
family, together with zero maps at the other coordinates, cannot
have a parameterization through $F_B(\tau_B)$.

\addcontentsline{toc}{section}{References}
\begin{thebibliography}{99}\small
\bibitem{BauerHanson} A.~Bauer and J.~E. Hanson, \emph{The Countable Reals}, \href{https://arxiv.org/abs/2404.01256v3}{arXiv:2404.01256v3}, 2026.
\bibitem{Bernardi} C.~Bernardi, \emph{Precomplete equivalence relations in dominical categories}, Proc. Amer. Math. Soc. \textbf{98} (1986), 124--134.
\bibitem{DPH} R.~A. Di Paola and A.~Heller, \emph{Dominical categories: recursion theory without elements}, J. Symbolic Logic \textbf{52} (1987), 594--635. \href{https://doi.org/10.2307/2274352}{doi:10.2307/2274352}.
\bibitem{Montagna} F.~Montagna, \emph{``Pathologies'' in two syntactic categories of partial maps}, Notre Dame J. Formal Logic \textbf{30} (1989), 105--116.
\bibitem{DPM} R.~A. Di Paola and F.~Montagna, \emph{Some properties of the syntactic $p$-recursion categories generated by consistent, recursively enumerable extensions of Peano arithmetic}, J. Symbolic Logic \textbf{56} (1991), 643--660.
\bibitem{CockettLack} J.~R.~B. Cockett and S.~Lack, \emph{Restriction categories. I. Categories of partial maps}, Theoret. Comput. Sci. \textbf{270} (2002), 223--259.
\bibitem{CockettHofstra} J.~R.~B. Cockett and P.~J.~W. Hofstra, \emph{Introduction to Turing categories}, Ann. Pure Appl. Logic \textbf{156} (2008), 183--209.
\bibitem{RangeII} J.~R.~B. Cockett, X.~Guo, and P.~Hofstra, \emph{Range categories II: Towards regularity}, Theory Appl. Categ. \textbf{26} (2012), 453--500. \href{https://www.tac.mta.ca/tac/volumes/26/18/26-18abs.html}{Journal text}.
\bibitem{Miller} J.~S. Miller, \emph{Degrees of unsolvability of continuous functions}, J. Symbolic Logic \textbf{69} (2004), 555--584. \href{https://doi.org/10.2178/jsl/1082418543}{doi:10.2178/jsl/1082418543}.
\bibitem{Keisler} H.~J. Keisler, \emph{The ultraproduct construction}, author's survey dated 1 June 2010, especially Theorem 3.1, Corollary 3.3, and Propositions 4.1--4.2. \href{https://people.math.wisc.edu/~hkeisler/ultraproducts-web-final.pdf}{Author's text}.
\bibitem{Terwijn} S.~Terwijn, \emph{The complexity of completions in partial combinatory algebra}, Math. Structures Comput. Sci. \textbf{34} (2024), 455--466, especially \S7. \href{https://doi.org/10.1017/S0960129524000252}{doi:10.1017/S0960129524000252}.
\bibitem{Aczel}
P.~Aczel, \emph{An introduction to inductive definitions},
in J.~Barwise (ed.), Handbook of Mathematical Logic,
North-Holland, 1977, pp.~739--782.
\href{https://lawrencecpaulson.github.io/papers/Aczel-Inductive-Defs.pdf}{Text}.
\bibitem{GTWW}
J.~A. Goguen, J.~W. Thatcher, E.~G. Wagner, and J.~B. Wright,
\emph{Initial algebra semantics and continuous algebras},
Journal of the ACM \textbf{24} (1977), no.~1, pp.~68--95.
\href{https://doi.org/10.1145/321992.321997}{doi:10.1145/321992.321997}.
\bibitem{CockettLists}
J.~R.~B. Cockett, \emph{List-arithmetic distributive categories: Locoi},
Journal of Pure and Applied Algebra \textbf{66} (1990), no.~1, pp.~1--29.
\href{https://doi.org/10.1016/0022-4049(90)90121-W}{doi:10.1016/0022-4049(90)90121-W}.
\bibitem{Wissmann}
T.~Wi\ss{}mann, \emph{Minimality notions via factorization systems and examples},
Logical Methods in Computer Science \textbf{18} (2022),
no.~3, article~31, pp.~31:1--31:26.
\href{https://doi.org/10.46298/lmcs-18(3:31)2022}{doi:10.46298/lmcs-18(3:31)2022}.
\end{thebibliography}
\end{document}
